\documentclass[12pt]{article}
\usepackage[margin=1in]{geometry}
\usepackage{amsmath}
\usepackage{graphicx}
\graphicspath{{./}{figures/}{plots/}{output/}}
\usepackage{booktabs}
\usepackage{setspace}
\usepackage{float}
\usepackage{caption}
\usepackage[round]{natbib}
\usepackage[colorlinks=true, linkcolor=black, citecolor=black, urlcolor=black]{hyperref}
\usepackage{xcolor}

\doublespacing

\title{Predictive Modeling of Ventricular Wall Stem Cell Fate During Postnatal Development}
\author{Jessica Yoon}
\date{}

\newcommand{\smartfig}[2][0.75\textwidth]{%
    \includegraphics[width=#1]{#2}%
}

\begin{document}

\maketitle

\begin{abstract}
Postnatal ventricular wall development involves the progressive transition from stem-cell-rich compartments toward a more mature ependymal state, but this process varies across anatomical regions and is difficult to observe continuously in vivo. In this study, I develop a compartment-specific stochastic model to predict stem cell composition across six ventricular wall regions defined by wall location (lateral versus medial) and rostrocaudal position (rostral, middle, caudal). The model is initialized using imaging-based measurements at postnatal day 7 (P7), including stem cell counts, ependymal cell counts, and observed division outcome probabilities. Development from P7 to P14 and P30 is simulated using a binomial model for the number of dividing stem cells and a multinomial model for division outcomes. I compare a baseline model without geometry to a crowding-aware model in which division intensity is modified by average apical area, a compartment-level proxy for local tissue geometry. Model performance is evaluated using predicted-versus-observed plots, residual analysis, mean absolute error (MAE), and root mean squared error (RMSE). The fitted crowding-aware model achieved an overall RMSE of 0.0815 and an overall MAE of 0.0603, suggesting that the framework captures broad developmental trends while also identifying compartments where prediction remains less reliable. Overall, this study provides an interpretable and biologically motivated approach for testing whether early fate behavior and local geometry can explain regional heterogeneity in ventricular wall maturation.
\end{abstract}

\section{Introduction}

Postnatal development of the ventricular wall depends on the coordinated balance between stem cell maintenance and differentiation into ependymal cells. Ependymal cells line the ventricular system and contribute to ventricular surface organization and cerebrospinal fluid circulation. Prior work has shown that ependymal cells are closely linked to radial glial lineages and that postnatal ventricular wall development involves both cellular fate specification and major changes in apical surface organization \citep{Spassky2005, Redmond2019}. Because of this, understanding how stem cells are depleted and replaced by mature ependymal cells is relevant not only to normal brain development but also to pathologic conditions in which ventricular lining integrity is disrupted.

A major challenge in studying this developmental process is that stem cell fate is not spatially uniform. The lateral and medial ventricular walls do not mature in the same way, and additional variation exists along the rostral, middle, and caudal axis. These regional differences suggest that a single global developmental rule is unlikely to explain all compartments equally well. At the same time, directly observing every division event across development is not feasible, especially when the available biological data consist of measurements at a small number of timepoints rather than continuous lineage tracing.

This project is motivated by the idea that even sparse imaging-based measurements can be used to build a predictive model if the major biological mechanisms are represented explicitly. In particular, stem cell division outcomes provide a natural probabilistic framework for modeling how a compartment changes over time. A dividing stem cell may undergo symmetric self-renewal, which increases the stem cell pool; asymmetric division, which preserves one stem cell while generating one ependymal cell; or symmetric differentiation, which generates two ependymal cells and reduces the stem cell pool. These outcomes are biologically meaningful because they determine whether a compartment remains stem-cell-rich or progresses more rapidly toward a mature ependymal state.

Another biologically relevant feature is tissue geometry. During postnatal development, the apical surface of ventricular wall cells changes substantially as ependymal cells mature and organize within the ventricular niche \citep{Redmond2019}. In this study, average apical area per cell is used as a compartment-level proxy for local crowding. Although apical area is not a direct measurement of mechanical force, contact inhibition, or cell-cell pressure, it provides a testable way to ask whether geometric information improves prediction of later stem cell composition.

The central goal of this study is to determine whether a simple compartment-specific stochastic model, initialized with P7 imaging data, can predict stem cell composition at P14 and P30 across six ventricular wall compartments. More specifically, this project asks two related questions. First, can a probabilistic model based on observed fate behavior reproduce broad developmental trends? Second, does incorporating a geometry-based crowding term improve prediction relative to a simpler baseline model without crowding? The first question evaluates whether early fate probabilities are sufficient to capture later composition. The second question evaluates whether local tissue geometry adds predictive value beyond fate probabilities alone. Framed this way, the project serves both as a biologically interpretable model of postnatal development and as a predictive modeling problem.

\section{Background and Related Work}

The ventricular wall is an important developmental niche because it contains cell populations involved in the transition from early neural progenitor states to more mature postnatal organization. Previous studies have shown that ependymal cells arise from radial glial populations and become largely postmitotic after differentiation \citep{Spassky2005}. More recent lineage and imaging work has emphasized that ependymal cells and postnatal neural stem cells can share common embryonic progenitors and that postnatal niche formation involves changes in apical area and spatial organization \citep{Redmond2019}. These findings support the biological structure of the current model: stem cell composition is expected to change over time, and spatial compartment may matter.

From a modeling perspective, cell fate decisions are often represented probabilistically because individual cells do not behave deterministically. Stochastic models are useful in developmental systems because they can represent uncertainty, heterogeneity, and the fact that population-level outcomes emerge from many individual fate events. For example, stochastic models of cell fate commitment have been used to simulate probabilistic transitions in differentiating cell populations and to connect observed biological measurements to underlying decision processes \citep{Teles2013}. The current study follows this general logic by separating the process into two interpretable steps: whether a stem cell divides, and what fate outcome occurs conditional on division.

This separation also motivates the statistical distributions used in the model. A binomial model is appropriate when a fixed number of cells each have a binary outcome, such as dividing or not dividing. A multinomial model extends this idea to more than two mutually exclusive outcomes, such as symmetric self-renewal, asymmetric division, and symmetric differentiation. These distributions are commonly used for categorical biological outcomes and provide a natural way to model cell-count changes while preserving the interpretation of each fate category \citep{HolmesHuber2019}.

\section{Data}

The analysis uses a merged imaging-based dataset, \texttt{combined\_modeling\_table.csv}, containing measurements from postnatal day 7 (P7), postnatal day 14 (P14), and postnatal day 30 (P30). Data are organized by compartment using two anatomical dimensions: wall location and regional position. Wall location is divided into lateral wall and medial wall, and each wall is further divided into rostral, middle, and caudal regions, yielding six biologically distinct compartments.

The dataset was generated from laboratory imaging measurements of the developing ventricular wall. Tissue samples were processed using the lab's experimental pipeline and summarized at the compartment level according to anatomical location. Cells were classified into stem cell and ependymal cell populations based on laboratory-defined criteria, and these classifications were then used to calculate compartment-specific counts, composition, and fate-related summaries. Because these data were generated within the lab rather than obtained from a public repository, the final submitted version should include the appropriate lab citation or acknowledgment according to advisor guidance. I acknowledge the Joanne Conover Lab for generating and sharing the ventricular wall imaging dataset used in this analysis.

Several types of measurements are included in the merged table. First, the data contain baseline counts of stem cells and ependymal cells at P7. These provide the initial conditions for model simulation. Second, the data include observed division outcome probabilities at P7, including probabilities for symmetric self-renewal, asymmetric division, and symmetric differentiation. Third, the dataset contains observed target composition at P14 and P30, including the fraction of stem cells and ependymal cells in each compartment. Finally, the data include average apical area per cell at later timepoints, which is used in the crowding-aware model as a proxy for local surface geometry.

The final analytic dataset is small. There are six P7 compartment-level initial conditions and twelve validation targets, corresponding to six compartments measured at two later timepoints (P14 and P30). This small sample size is important for interpretation because each compartment contributes only two observed target values to the validation summary. As a result, compartment-specific error estimates are useful for identifying possible patterns, but they should not be interpreted as definitive evidence that one region is biologically easier or harder to predict.

Before simulation and validation, the merged table was screened for missing values, inconsistent compartment labels, negative counts, impossible proportions, and division outcome probabilities that did not sum to one within a compartment. Rows with missing model inputs were excluded from the final analysis. These preprocessing checks are especially important because the modeling table combines multiple imaging-derived summaries into a single analytical dataset.

\section{Methods}

\subsection{Overview of modeling strategy}

The model tracks two cell populations within each compartment: the number of stem cells, denoted by $S(t)$, and the number of ependymal cells, denoted by $E(t)$. The system is initialized at P7 using observed counts from the imaging data. Development is then modeled in coarse steps from P7 to P14 and from P14 to P30. The primary outcome of interest is the predicted stem cell fraction,
\[
\text{SC fraction}(t) = \frac{S(t)}{S(t)+E(t)}.
\]

The analysis includes two related models. The first is a baseline model in which the division fraction is fixed within each compartment and does not depend on geometry. The second is a crowding-aware model in which the division fraction depends on average apical area through a sigmoid function. Comparing these two models allows the analysis to test whether including a geometry-based predictor improves predictive accuracy.

\subsection{Model assumptions}

The model relies on several simplifying assumptions. First, ependymal cells are treated as non-dividing over the modeled intervals, so all changes in compartment composition arise from stem cell division outcomes. Second, within a given compartment and time step, cells are assumed to share the same division fraction and the same fate outcome probabilities. Third, fate probabilities estimated at P7 are used as proxies for later developmental transitions. Fourth, average apical area per cell is treated as a compartment-level proxy for local crowding rather than a direct mechanistic measurement of physical constraint. Fifth, the model does not explicitly include cell death, migration between compartments, or measurement error in the observed counts.

These assumptions are useful because they keep the model interpretable, but they also define the limits of the analysis. The most biologically restrictive assumption is that fate probabilities remain constant after P7. If fate behavior changes substantially between P7, P14, and P30, then the model may incorrectly attribute later compositional change to initial fate probabilities or crowding effects. Similarly, ignoring cell death, migration, and measurement error means that all observed changes are explained through division and fate assignment. This could bias compartment-specific predictions if some regions are affected by processes not included in the model. Therefore, the model should be interpreted as a simplified predictive approximation rather than a complete mechanistic description of ventricular wall development.

\subsection{Division model}

At each modeled time step, only a subset of stem cells is assumed to divide. Let $S(t)$ denote the number of stem cells in a compartment at time $t$, and let $r(t)$ denote the fraction of stem cells expected to divide during that interval. The number of dividing stem cells is modeled as
\[
N_{\mathrm{div}}(t) \sim \mathrm{Binomial}(S(t), r(t)).
\]

A binomial distribution is appropriate here because each stem cell is treated as having two possible states within a time interval: divide or do not divide. Under this formulation, the total number of division events depends on the size of the stem cell pool and the probability of division.

\subsection{Fate assignment model}

Conditional on division, each dividing stem cell is assigned one of three mutually exclusive outcomes: symmetric self-renewal, asymmetric division, or symmetric differentiation. Let $p_{SS}$, $p_{SE}$, and $p_{EE}$ denote the probabilities of these three outcomes, respectively, with
\[
p_{SS}+p_{SE}+p_{EE}=1.
\]

Then the vector of fate outcomes is modeled as
\[
(n_{SS}, n_{SE}, n_{EE}) \sim \mathrm{Multinomial}(N_{\mathrm{div}}(t), p_{SS}, p_{SE}, p_{EE}).
\]

A multinomial distribution is appropriate because each dividing stem cell must fall into exactly one of three mutually exclusive fate categories. This separation of the model into a binomial division step and a multinomial fate step mirrors the biological process: first a stem cell either divides or does not divide, and then a dividing cell takes one of several possible outcomes.

\subsection{Population update equations}

After the number of divisions and their outcomes are sampled, the compartment populations are updated according to
\[
S(t+1)=S(t)+n_{SS}-n_{EE}
\]
and
\[
E(t+1)=E(t)+n_{SE}+2n_{EE}.
\]

These equations reflect the biological meaning of each fate. Symmetric self-renewal increases the stem cell pool, asymmetric division preserves the stem cell pool while adding one ependymal cell, and symmetric differentiation reduces the stem cell pool while adding two ependymal cells.

\subsection{Baseline model}

In the baseline model, the division fraction $r(t)$ is fixed within each compartment and does not depend on average apical area. All other components of the simulation remain the same, including the binomial division step, multinomial fate assignment, population update equations, and repeated Monte Carlo simulation. This model serves as the reference point for evaluating whether the addition of a crowding-related geometry term improves predictive performance. In other words, the baseline model asks whether P7 counts and P7 fate probabilities alone are enough to reproduce later stem cell composition.

\subsection{Crowding-aware model}

In the crowding-aware model, the division fraction is allowed to vary as a function of average apical area. The rationale is that local geometry may influence how actively stem cells continue to divide as the ventricular wall matures. This relationship is modeled using a sigmoid function,
\[
r(t)=\frac{1}{1+e^{-(\alpha-\beta a(t))}},
\]
where $\alpha$ and $\beta$ are crowding parameters and $a(t)$ denotes the average apical area for a compartment at time $t$.

The sigmoid function was chosen because it keeps the division fraction bounded between 0 and 1, which is necessary for a probability-like quantity. The parameter $\alpha$ controls the baseline tendency toward division when the area term is small, while $\beta$ controls how strongly increasing apical area reduces the expected division fraction. Thus, larger values of $\alpha$ increase division intensity, while larger values of $\beta$ make the model more sensitive to crowding.

In the current implementation, $\alpha$ and $\beta$ were treated as fixed hyperparameters rather than freely estimated from the data. This choice was made because the validation dataset contains only twelve target observations, so freely estimating multiple crowding parameters would create a high risk of overfitting and poor identifiability. Fixing $\alpha$ and $\beta$ allows the crowding term to represent a biologically motivated sensitivity assumption while keeping the model simple. However, this decision also means that the final values must be reported clearly for reproducibility.

In the final crowding-aware model implementation, the crowding parameters were fixed at $\alpha = 0.25$ and $\beta = 0.80$. These parameters were not estimated freely from the data because the validation dataset was small, with only six compartments evaluated at two postnatal timepoints. Instead, they were treated as fixed hyperparameters chosen to encode a moderate, biologically plausible relationship between apical area and division activity while avoiding overfitting.

In the code, apical area was first standardized as
\[
z = \frac{a - \bar{a}}{s_a},
\]
where $a$ is the compartment-specific apical area, $\bar{a}$ is the mean apical area, and $s_a$ is the standard deviation of apical area. The division fraction was then calculated as
\[
r(t)=\frac{1}{1+e^{-(0.25-0.80z)}}.
\]
Under this parameterization, compartments with larger-than-average apical area have lower predicted division activity, while compartments with smaller-than-average apical area have higher predicted division activity. The value $\alpha=0.25$ sets the baseline division tendency when apical area is near the mean, and $\beta=0.80$ controls how strongly division activity changes as apical area moves above or below the mean. This choice allowed the model to include a geometry-informed crowding effect while keeping the relationship simple, bounded between 0 and 1, and biologically interpretable.

\subsection{Model calibration}

Rather than fitting a large number of free parameters, the final model used fixed, interpretable calibration choices to reduce overfitting in the small validation dataset. The baseline division fraction was set to $r=0.30$. In the crowding-aware model, apical area modified this division fraction using fixed hyperparameters $\alpha=0.25$ and $\beta=0.80$. These values were chosen to encode a moderate negative relationship between standardized apical area and division activity while keeping predicted division fractions bounded between 0 and 1.

Model performance was then evaluated by comparing predicted and observed stem cell fractions at P14 and P30. The baseline and crowding-aware models were compared using the same validation targets and the same error metrics, including MAE and RMSE. This approach emphasizes predictive comparison rather than formal parameter estimation, which is appropriate given the limited number of compartment-timepoint observations.

\subsection{Monte Carlo simulation}

Because stem cell fate is probabilistic, the model is simulated repeatedly rather than run only once. For each compartment, many stochastic trajectories are generated from the same initial conditions and model parameters. These repeated simulations produce a predictive distribution for the stem cell fraction at P14 and P30 rather than a single deterministic value.

In the current implementation, compartment trajectories were simulated across the available developmental transitions and summarized over 3000 Monte Carlo prediction runs per model.

\subsection{Validation}

Model performance is evaluated by comparing predicted stem cell composition at P14 and P30 with observed values from the dataset. Validation includes both graphical and quantitative components. Predicted-versus-observed scatterplots provide an overall visual summary of fit across compartment-timepoint combinations. Residual plots are used to assess whether the model systematically underpredicts or overpredicts at specific timepoints. Model performance is summarized numerically using both mean absolute error (MAE) and root mean squared error (RMSE).

The primary quantitative validation metric is RMSE, defined as
\[
\mathrm{RMSE}=\sqrt{\frac{1}{n}\sum_{i=1}^{n}(\hat{y}_i-y_i)^2},
\]
where $\hat{y}_i$ is the predicted stem cell fraction and $y_i$ is the observed stem cell fraction. RMSE is useful here because it summarizes overall prediction error while placing greater weight on larger deviations, which helps identify compartments where the model performs poorly. MAE is also reported as a complementary metric because it is easier to interpret directly as the average absolute discrepancy between predicted and observed values.

Because the validation set contains only twelve target observations, these metrics should be interpreted cautiously. An RMSE of 0.0815 means that the model's typical prediction error is approximately seven percentage points of stem cell fraction. In this context, that level of error suggests reasonable recovery of broad developmental trends, but it is not strong enough to support highly precise compartment-level conclusions.

\section{Results}

\subsection{Summary of baseline compartment differences}

Preliminary inspection of the data suggests that the six ventricular wall compartments begin with meaningfully different baseline conditions. P7 stem cell counts vary across compartments, indicating that the initial stem cell pool is not uniform throughout the ventricular wall. Likewise, the observed division outcome probabilities differ by compartment, supporting the biological premise that developmental behavior is region-specific rather than globally constant.

These baseline differences are important because they directly influence the simulations. Compartments with larger starting stem cell pools tend to produce smoother predicted trajectories because stochastic variation is averaged over more cells. Compartments with fewer initial stem cells are more sensitive to random variation, so a small number of division events can meaningfully shift the predicted stem cell fraction. Similarly, compartments with more strongly differentiation-weighted fate probabilities are expected to show faster depletion of the stem cell pool over time.

Figure 1 summarizes P7 stem cell counts by compartment, and Figure 2 displays the fate probability composition for each compartment.

\begin{figure}[H]
\centering
\smartfig{p7_counts.png}
\caption{P7 stem cell counts by compartment.}
\label{fig:p7_counts}
\end{figure}

\begin{figure}[H]
\centering
\smartfig{fate_probabilities.png}
\caption{Compartment-specific division outcome probabilities at P7.}
\label{fig:fate_probs}
\end{figure}

\subsection{Observed developmental targets}

Observed stem cell composition at P14 and P30 indicates that postnatal ventricular wall development is characterized by progressive depletion of the stem cell population, although the magnitude and timing of this depletion differ across compartments. This supports the idea that lateral and medial wall regions do not mature in identical ways.

Average apical area also varies across compartments and over time. When compared with observed stem cell composition, the geometry proxy appears to capture biologically relevant variation, suggesting that physical arrangement or packing may be associated with how quickly a compartment transitions toward a more mature ependymal state. However, this relationship should be interpreted as associative rather than causal because average apical area is a compartment-level summary and not a direct measurement of mechanical crowding.

\begin{figure}[H]
\centering
\smartfig{area_vs_sc.png}
\caption{Average apical area versus observed stem cell fraction.}
\label{fig:area_vs_sc}
\end{figure}

\subsection{Baseline and crowding-aware model comparison}

The baseline model provides a useful reference point for evaluating whether fate probabilities alone are sufficient to capture compartment-specific developmental change. The crowding-aware model tests whether adding average apical area improves prediction beyond that baseline. This comparison is important because it turns the crowding term into an evaluable modeling choice rather than simply a descriptive biological idea.

Table~\ref{tab:model_comparison} compares both models using the same validation targets and the same error metrics. The baseline model represents predictions based on P7 counts and fate probabilities alone, while the crowding-aware model additionally incorporates average apical area as a geometry-based predictor.

\begin{table}[H]
\centering
\caption{Overall model comparison using the same validation targets.}
\label{tab:model_comparison}
\begin{tabular}{lccc}
\toprule
Model & Overall MAE & Overall RMSE & Interpretation \\
\midrule
Baseline, no crowding & 0.0808 & 0.0976 & Fate probabilities only \\
Crowding-aware & 0.0603 & 0.0815 & Fate probabilities + apical area \\
\bottomrule
\end{tabular}
\end{table}

Compared with the baseline model, the crowding-aware model modestly improved overall predictive performance. The baseline model had an overall MAE of 0.0808 and an overall RMSE of 0.0976, while the crowding-aware model had an overall MAE of 0.0603 and an overall RMSE of 0.0815. This suggests that adding the apical-area-based crowding term improved prediction accuracy, although the improvement should be interpreted cautiously because the validation set is small.

\subsection{Crowding-aware model performance}

The crowding-aware model achieved an overall RMSE of \textbf{0.0815} and an overall MAE of \textbf{0.0603}. Interpreted on the stem cell fraction scale, this means that the model's average prediction error was on the order of approximately six to eight percentage points. This is a reasonable level of error for capturing broad developmental trends, especially given the small number of validation observations, but it should not be interpreted as evidence of highly precise compartment-level prediction.

Across compartments, prediction error was lowest in the \textbf{medial caudal} compartment (RMSE $= 0.0276$, MAE $= 0.0262$) and highest in the \textbf{lateral caudal} compartment (RMSE $= 0.1176$, MAE $= 0.1130$). More broadly, the medial wall compartments showed lower error than the lateral wall compartments in this fitted version of the model. One possible biological interpretation is that medial wall development may be better captured by the current fate and crowding assumptions. A more cautious statistical interpretation is that lateral compartments may have greater variability, smaller effective counts, or developmental dynamics not represented in the current model. Because each compartment contributes only two evaluation points, these regional patterns should be treated as suggestive rather than definitive.

\begin{table}[H]
\centering
\caption{Compartment-specific prediction error for the crowding-aware model.}
\label{tab:compartment_errors}
\begin{tabular}{lccc}
\toprule
Compartment & MAE & RMSE & n \\
\midrule
Medial / Caudal   & 0.0262 & 0.0276 & 2 \\
Medial / Rostral  & 0.0319 & 0.0354 & 2 \\
Medial / Middle   & 0.0332 & 0.0409 & 2 \\
Lateral / Rostral & 0.0705 & 0.0779 & 2 \\
Lateral / Middle  & 0.0742 & 0.0861 & 2 \\
Lateral / Caudal  & 0.1130 & 0.1176 & 2 \\
\bottomrule
\end{tabular}
\end{table}

Table~\ref{tab:compartment_errors} shows that model performance was not uniform across compartments. The lateral caudal compartment had the poorest fit overall, which may indicate that this region has biological behavior not captured by the current model assumptions. For example, its developmental trajectory may depend on time-varying fate probabilities, additional spatial cues, or measurement variability that is not represented by average apical area alone. However, because each compartment-specific metric is based on only two points, these values should be used mainly to guide future model refinement rather than to make strong claims about regional biology.

\subsection{Predicted-versus-observed comparisons}

Predicted-versus-observed comparisons indicate that the model captures the overall developmental direction of stem cell depletion, but fit varies meaningfully across compartments. Points close to the diagonal line represent accurate predictions, while points farther from the diagonal show where the model either underpredicts or overpredicts stem cell fraction. The residual plot provides a complementary view by showing whether errors cluster by compartment or timepoint.

\begin{figure}[H]
\centering
\smartfig{predicted_vs_observed.png}
\caption{Predicted-versus-observed stem cell fraction.}
\label{fig:pred_obs}
\end{figure}

\begin{figure}[H]
\centering
\smartfig{residual_plot.png}
\caption{Residual plot by compartment or timepoint.}
\label{fig:residuals}
\end{figure}

Together, these plots suggest that the model is more successful at recovering the general direction of postnatal stem cell depletion than at perfectly matching every compartment-timepoint combination. This is consistent with the purpose of the model: it is intended to test whether a simple, interpretable stochastic process can explain broad developmental structure, not to provide a fully mechanistic account of every regional difference.

\section{Discussion}

This study presents a compartment-specific probabilistic framework for modeling postnatal ventricular wall development from sparse imaging-based data. One major strength of the approach is that it remains biologically interpretable throughout. Rather than predicting later cell composition from arbitrary features alone, the model explicitly encodes developmental mechanisms in terms of stem cell division, fate choice, compartment-specific initial conditions, and an optional geometry-linked crowding effect.

A second strength is that the framework makes uncertainty visible. Because the model is simulated many times, it produces predictive distributions rather than single deterministic estimates. This is especially important in a biological setting where small cell populations and probabilistic fate decisions can generate substantial variability. The use of Monte Carlo simulation is therefore not just a computational choice; it matches the biological idea that repeated developmental trajectories could differ even under the same starting conditions.

The comparison between the baseline and crowding-aware models is also valuable from a data science perspective. A baseline model asks whether initial counts and fate probabilities alone can explain later stem cell composition. The crowding-aware model asks whether adding average apical area provides additional predictive information. This structure is important because it prevents the analysis from assuming that the crowding term is useful simply because it is biologically plausible. Instead, the crowding term should be judged by whether it improves prediction or provides a clearer interpretation relative to the simpler model.

The fitted crowding-aware model achieved an overall RMSE of 0.0815 and an overall MAE of 0.0603. These values suggest that the model can recover broad developmental trends on the stem cell fraction scale, but they should be interpreted in light of the small validation sample. An error of roughly six to eight percentage points is promising for a simple model, but the dataset is not large enough to support strong claims about precise regional ranking. The lower error in medial compartments and higher error in the lateral caudal compartment may reflect meaningful biological differences, but it could also reflect greater noise, smaller effective sample size, or omitted mechanisms in the lateral wall.

Several limitations should be acknowledged. First, ependymal cells are assumed not to divide over the modeled intervals. Second, the model does not explicitly account for cell death, migration, or measurement uncertainty. Third, development is represented using coarse transitions rather than a continuous-time process. Fourth, fate probabilities are held constant from P7 onward, even though actual fate behavior may change during development. Fifth, the crowding parameters $\alpha$ and $\beta$ are fixed hyperparameters rather than fully estimated from the data. This improves simplicity and reduces overfitting risk, but it also means that conclusions about crowding depend on the chosen parameter values.

These simplifying assumptions may affect prediction error in different ways. For example, if ependymal cells divide at low frequency, then the current model may understate later ependymal expansion. If stem cell fate probabilities shift substantially after P7, then using fixed early probabilities may oversimplify later developmental transitions. Ignoring cell death or migration may also attribute all compositional change to division dynamics, which could bias the interpretation of compartment-specific effects. For this reason, the current framework should be interpreted as a simplified predictive approximation rather than a complete mechanistic description of ventricular wall development.

These limitations do not make the model uninformative. Rather, they define the scope of interpretation. At this stage, the model is most useful as a proof of concept showing that sparse compartment-level imaging data can be translated into an interpretable stochastic simulation. It also identifies where the model performs less well, which is valuable for future refinement. In particular, the higher errors observed in the lateral wall compartments suggest that these regions may require additional biological predictors, time-varying fate probabilities, or a more flexible compartment-specific structure.

There are several directions for future refinement. A future improvement would be to expand this side-by-side comparison with additional validation data and sensitivity analyses, rather than relying on only twelve validation targets. A second improvement would be to run a sensitivity analysis over different fixed values of $\alpha$ and $\beta$ to determine whether the crowding-aware results are stable. A third improvement would be to allow fate probabilities to vary over time instead of relying only on P7 estimates. Additional validation metrics, larger datasets, replicate-level modeling, and partial pooling across compartments would also strengthen the final assessment of compartment-specific fit.

\section{Conclusion}

This study develops a predictive model of ventricular wall stem cell fate during postnatal development using compartment-specific imaging data and stochastic simulation. By combining baseline cell counts, observed division outcome probabilities, and a geometry-based crowding proxy, the model predicts stem cell composition at later timepoints and compares those predictions with observed data.

The fitted crowding-aware model achieved an overall RMSE of 0.0815 and an overall MAE of 0.0603, suggesting that the model captures broad developmental trends while also revealing regional differences in fit. However, because the validation dataset contains only twelve target observations, compartment-specific differences should be interpreted cautiously. The main contribution of this project is therefore not that it definitively identifies one compartment as easier or harder to model, but that it provides a coherent and biologically interpretable framework for evaluating how early fate behavior and tissue geometry may shape postnatal ventricular wall development. Future work with larger datasets, stronger baseline comparisons, and sensitivity analysis of the crowding parameters would make the model more reproducible and biologically informative.

\begin{thebibliography}{9}

\bibitem[Holmes and Huber(2019)]{HolmesHuber2019}
Holmes, S., \& Huber, W. (2019).
\textit{Modern statistics for modern biology}.
Cambridge University Press.

\bibitem[Redmond et~al.(2019)]{Redmond2019}
Redmond, S. A., Figueres-O\~{n}ate, M., Obernier, K., Nascimento, M. A., Parraguez, J. I., L\'{o}pez-Mascaraque, L., Fuentealba, L. C., \& Alvarez-Buylla, A. (2019).
Development of ependymal and postnatal neural stem cells and their origin from a common embryonic progenitor.
\textit{Cell Reports, 27}(2), 429--441.e3.

\bibitem[Spassky et~al.(2005)]{Spassky2005}
Spassky, N., Merkle, F. T., Flames, N., Tramontin, A. D., Garc\'{i}a-Verdugo, J. M., \& Alvarez-Buylla, A. (2005).
Adult ependymal cells are postmitotic and are derived from radial glial cells during embryogenesis.
\textit{Journal of Neuroscience, 25}(1), 10--18.

\bibitem[Teles et~al.(2013)]{Teles2013}
Teles, J., Pina, C., Eden, P., Ohlsson, M., Enver, T., \& Peterson, C. (2013).
Transcriptional regulation of lineage commitment: A stochastic model of cell fate decisions.
\textit{PLOS Computational Biology, 9}(8), e1003197.

\end{thebibliography}

\end{document}
